\subsection{Convergence of the relational weight construction in the dense graph limit}
  \label{subsec:convergence_dense}

  The relational weight construction is defined for finite graphs.
  We examine its compatibility with the dense limit $N\to\infty$, where the neighbourhood
  radius $\ell_0(N)$ is scaled appropriately.

  The combinatorial neighbourhood $\mathcal{N}_i$ of radius $\ell_0$
  corresponds to a ball of radius $O(\ell_0/N^{1/d})$ in the ambient
  metric as the graph becomes dense, with $d$ the intrinsic spectral
  dimension.
  Under standard bandwidth conditions on $\ell_0(N)$ --- namely
  $\ell_0(N)\to 0$ and $N\,\ell_0(N)^d \to \infty$ for i.i.d.\ sampling
  from a $C^2$ density on a compact manifold --- the local neighbourhood
  aggregates $\bar{\chi}_i$ defined in Eq.~\eqref{eq:chi-bar} converge
  pointwise to a smooth function on the underlying manifold.

  The weight functional~\eqref{eq:weights} satisfies the smoothness and
  locality assumptions required by the convergence results of
  Belkin and Niyogi~\cite{Belkin2003} and Hein, Audibert, and von
  Luxburg~\cite{Hein2007}.
  Consequently, the weighted graph Laplacian converges to the
  Laplace\textendash Beltrami operator (modulated by the local relational
  stiffness), the effective spectral description converges to a smooth
  coarse-grained representation, and the spectral distances approximate
  the corresponding continuous diffusion distances.

  Thus, under standard regularity assumptions, the relational weight
  construction is consistent with the continuum limit.
